\documentclass[11pt,a4paper]{article}
\usepackage[T1]{fontenc}
\usepackage[utf8]{inputenc}
\usepackage{lmodern}
\usepackage{amsmath,amssymb}
\usepackage{booktabs}
\usepackage{array}
\usepackage[a4paper,margin=2.2cm]{geometry}
\title{Minimal QUBO for Telecom Backhaul Allocation}
\author{}
\date{}

\begin{document}
\maketitle

\textbf{Goal.} Allocate limited backhaul bandwidth to service streams while respecting SLAs.
The problem is cast as a Quadratic Unconstrained Binary Optimisation (QUBO) so it can be
solved directly on a D-Wave quantum annealer or any compatible sampler.

% -----------------------------------------------------------------------
\section*{Minimal Case}

Two streams $s\in\{1,2\}$ share one bottleneck link with capacity $C$ Mbps.
The solver must decide, in discrete \emph{chunk} steps of size $\Delta$ Mbps, how much
bandwidth each stream receives.

% -----------------------------------------------------------------------
\section*{Variable Glossary}

\subsection*{Problem Parameters (inputs, not optimised)}

\begin{tabular}{@{}lp{10.5cm}@{}}
\toprule
Symbol & Meaning \\
\midrule
$C$ & Link capacity in Mbps. Hard upper bound on the total bandwidth that can be
      assigned across all streams sharing that link.
      \textit{Code:} \texttt{Link.capacity\_mbps}. \\[4pt]

$\Delta$ & Chunk size in Mbps. The smallest unit of bandwidth that can be
           allocated. A finer $\Delta$ gives more resolution but adds more binary
           variables.
           \textit{Code:} \texttt{delta\_mbps} (default 10). \\[4pt]

$D_s$ & Demand of stream $s$ in Mbps: the bandwidth it ideally wants.
        \textit{Code:} \texttt{Stream.demand\_mbps}. \\[4pt]

$G_s$ & Minimum guarantee of stream $s$ in Mbps: the floor below which the
        allocation must not fall (SLA floor).
        \textit{Code:} \texttt{Stream.min\_guarantee\_mbps}. \\[4pt]

$w_s$ & Priority weight of stream $s$. A higher value means shortfall for this
        stream is penalised more heavily. Voice streams typically carry
        $w_s \approx 3$--$4$; best-effort data carries $w_s \approx 1$.
        \textit{Code:} \texttt{Stream.priority\_weight}. \\[4pt]

$B_s^{\text{prev}}$ & Bandwidth allocated to stream $s$ in the \emph{previous}
                      scheduling interval. Used to penalise large swings that
                      would disrupt in-flight flows.
                      \textit{Code:} \texttt{Stream.prev\_alloc\_mbps}. \\[4pt]

$K_s$ & Number of allocation chunks available for stream $s$,
        $K_s = \lceil \max(D_s, G_s) / \Delta \rceil$.
        \textit{Code:} \texttt{max\_chunks\_by\_stream[s.name]}. \\[4pt]

$\lambda_{\text{sf}}$ & Penalty strength for the shortfall equality constraint.
                        Governs how strongly the solver is pushed to satisfy
                        $B_s + q_s = D_s$.
                        \textit{Code:} \texttt{lambda\_shortfall} (default 6.0). \\[4pt]

$\lambda_{\text{st}}$ & Penalty strength for the stability term.  Kept small
                        relative to $\lambda_{\text{sf}}$ so that stability is a
                        soft preference, not a hard requirement.
                        \textit{Code:} \texttt{lambda\_stability} (default 0.03). \\[4pt]

$\lambda_{\text{cap}}$ & Penalty strength for the link-capacity constraint.
                         Should be the largest $\lambda$ to make capacity
                         violation the most expensive outcome.
                         \textit{Code:} \texttt{lambda\_capacity} (default 15.0). \\[4pt]

$\lambda_{\text{min}}$ & Penalty strength for the minimum-guarantee constraint.
                         Set between $\lambda_{\text{sf}}$ and $\lambda_{\text{cap}}$
                         to prioritise guarantees over full-demand satisfaction.
                         \textit{Code:} \texttt{lambda\_min\_guarantee} (default 12.0). \\[4pt]

$\alpha$ & Linear cost coefficient applied to the shortfall.
           Together with $w_s$ it provides a \emph{direct} per-Mbps shortfall
           cost in addition to the quadratic penalty.
           \textit{Code:} embedded as \texttt{sf\_linear\_weight = lambda\_shortfall * s.priority\_weight}. \\

\bottomrule
\end{tabular}

\subsection*{Decision Variables (binary, optimised by the solver)}

\begin{tabular}{@{}lp{10.5cm}@{}}
\toprule
Symbol & Meaning \\
\midrule
$y_{s,k}\in\{0,1\}$ & Allocation bit for stream $s$, chunk index $k=0,\dots,K_s-1$.
                       Setting $y_{s,k}=1$ grants $\Delta$ Mbps to stream $s$.
                       \textit{Code variable name:} \texttt{y::\{stream\}::\{k\}}. \\[4pt]

$B_s$ & Derived (not a QUBO variable itself): total bandwidth allocated to
        stream $s$,
        \[B_s = \Delta\sum_{k=0}^{K_s-1} y_{s,k}.\]
        It is a linear function of the $y_{s,k}$ bits and appears inside every
        penalty term. \\[4pt]

$q_{s,b}\in\{0,1\}$ & Shortfall bit $b$ for stream $s$.
                       Together these bits binary-encode the integer shortfall
                       in chunk units:
                       \[q_s = \Delta\sum_{b} 2^b\, q_{s,b}.\]
                       $q_s$ represents how many Mbps stream $s$ falls short of
                       its demand $D_s$.  It is introduced so that the equality
                       $B_s + q_s = D_s$ can be written as a quadratic penalty.
                       \textit{Code variable name:} \texttt{sf::\{stream\}::\{b\}}.
                       Number of bits: $\lceil\log_2 K_s\rceil + 1$.\\[4pt]

$u_b\in\{0,1\}$ & Capacity slack bit $b$ for the link.
                  Together these bits binary-encode the unused link capacity:
                  \[u = \Delta\sum_{b} 2^b\, u_b,\quad u \ge 0.\]
                  $u$ lets the sum $B_1+B_2+u=C$ be turned into a
                  penalty that is zero exactly when the constraint is tight or
                  under-utilised.
                  \textit{Code variable name:} \texttt{csl::\{link\}::\{b\}}.\\[4pt]

$r_{s,b}\in\{0,1\}$ & Guarantee-slack bit $b$ for stream $s$.
                       Together these bits binary-encode the surplus above the
                       guarantee floor:
                       \[r_s = \Delta\sum_{b} 2^b\, r_{s,b},\quad r_s \ge 0.\]
                       $r_s$ captures how far above $G_s$ the allocation sits.
                       The equality $B_s - r_s = G_s$ is zero only when
                       $B_s \ge G_s$, so any violation raises the energy.
                       \textit{Code variable name:} \texttt{gsl::\{stream\}::\{b\}}.\\

\bottomrule
\end{tabular}

\bigskip
\noindent\textbf{Why binary encode continuous slack?}
QUBO requires every variable to be binary ($\{0,1\}$).  Non-negative integers
like $q_s$ are represented in binary via the standard basis
$q_s = \sum_b 2^b q_{s,b}$, scaled by $\Delta$ to stay in Mbps units.
The number of bits needed is $\lceil\log_2(\text{max value})\rceil + 1$, computed
in the code by \texttt{bit\_width\_for\_nonnegative\_integer}.

% -----------------------------------------------------------------------
\section*{Objective Terms}

We minimise a weighted sum of penalty terms.  Each term is zero when the
corresponding constraint or preference is exactly satisfied, and grows
quadratically with the violation.

\subsection*{1) Shortfall (soft demand satisfaction)}

We want $B_s$ to be as close to $D_s$ as possible.  Introduce the shortfall
$q_s \ge 0$ and enforce the equality:
\[
B_s + q_s = D_s,\quad q_s\ge 0.
\]
The quadratic penalty drives the equality to zero; the linear term gives a
per-Mbps cost proportional to the stream's priority weight:
\[
\lambda_{\text{sf}}\sum_{s=1}^2 (B_s + q_s - D_s)^2
+ \alpha\sum_{s=1}^2 w_s\, q_s.
\]
When allocation equals demand ($B_s = D_s$) the optimal $q_s = 0$ and both
terms vanish.  When the link is over-subscribed, the solver raises $q_s$ for
lower-priority streams first because their $w_s$ is smaller, incurring a
lower linear cost.

\subsection*{2) Stability (allocation smoothness)}

Large shifts between intervals disrupt TCP flows and increase re-ordering.
We penalise the squared distance from the previous allocation:
\[
\lambda_{\text{st}}\sum_{s=1}^2 (B_s - B_s^{\text{prev}})^2.
\]
$\lambda_{\text{st}} \ll \lambda_{\text{sf}}$ so this is a soft preference:
the solver will deviate from $B_s^{\text{prev}}$ when necessary to satisfy
demand or capacity, but avoids unnecessary churn.

% -----------------------------------------------------------------------
\section*{Constraints as Quadratic Penalties}

Hard constraints cannot be expressed in QUBO directly; instead they are added
as penalty terms whose coefficients are large enough that any constraint
violation costs more than any objective gain.

\subsection*{3) Link capacity}

The total bandwidth on the link must not exceed $C$.  Introduce nonnegative
slack $u$ representing unused capacity:
\[
B_1 + B_2 + u = C,\quad u\ge 0,
\]
\[
\lambda_{\text{cap}}(B_1 + B_2 + u - C)^2.
\]
$u$ absorbs the gap between total allocation and the capacity ceiling.  When the
link is exactly full, $u=0$.  When under-utilised, $u=C - B_1 - B_2 > 0$.
Over-subscription ($B_1+B_2>C$) cannot be absorbed by $u$ (it is non-negative),
so the penalty is non-zero, repelling the solver from that region.
$\lambda_{\text{cap}}$ is the largest penalty coefficient to make capacity
violation the most expensive outcome in the energy landscape.

\subsection*{4) Minimum guarantees}

Each stream must receive at least $G_s$ Mbps.  A nonnegative slack $r_s$
captures the surplus:
\[
B_s - r_s = G_s,\quad r_s\ge 0,
\]
\[
\lambda_{\text{min}}\sum_{s=1}^2 (B_s - r_s - G_s)^2.
\]
When $B_s \ge G_s$ the solver can set $r_s = B_s - G_s \ge 0$ and the penalty
vanishes.  When $B_s < G_s$ there is no valid non-negative $r_s$; the penalty
is strictly positive, steering the solver away from under-provisioning.

% -----------------------------------------------------------------------
\section*{Final QUBO Energy}

Combining all terms gives the energy function that the annealer minimises:
\begin{align*}
E(\mathbf{x}) &=
\underbrace{\lambda_{\text{sf}}\sum_{s=1}^2 (B_s + q_s - D_s)^2
+ \alpha\sum_{s=1}^2 w_s q_s}_{\text{shortfall cost}} \\[6pt]
&\quad+\;\underbrace{\lambda_{\text{st}}\sum_{s=1}^2 (B_s - B_s^{\text{prev}})^2}_{\text{stability cost}}
+\;\underbrace{\lambda_{\text{cap}}(B_1 + B_2 + u - C)^2}_{\text{capacity constraint}}
+\;\underbrace{\lambda_{\text{min}}\sum_{s=1}^2 (B_s - r_s - G_s)^2}_{\text{guarantee constraint}}.
\end{align*}

\noindent All variables are binary after encoding $q_s$, $u$, $r_s$ with bit vectors.
The expansion of each squared term produces the linear (diagonal) and quadratic
(off-diagonal) QUBO coefficients added by \texttt{add\_square\_penalty} in the
code.

\bigskip
\noindent\textbf{Penalty coefficient ordering.} For constraints to dominate
objectives the coefficients must satisfy
\[
\lambda_{\text{cap}} > \lambda_{\text{min}} > \lambda_{\text{sf}} \gg \lambda_{\text{st}}.
\]
The defaults (15, 12, 6, 0.03) reflect this hierarchy: capacity is never
violated, guarantees are protected next, demand satisfaction comes third, and
stability is a soft tie-breaker.

% -----------------------------------------------------------------------
\section*{Tiny Numeric Example}

Let $\Delta=10$ Mbps, $C=100$, $D_1=70$, $D_2=60$, $G_1=40$, $G_2=30$.
Then $D_1+D_2=130>C$, so the optimizer must throttle.

Stream 1 and 2 each get $K_s = \lceil D_s/\Delta \rceil = 7$ and $6$ allocation
bits respectively, plus shortfall and slack bits, for a total variable count
well within classical sampler reach.

If $w_1 > w_2$, the shortfall linear term penalises stream 2's shortfall less
than stream 1's.  The optimal QUBO solution therefore keeps $B_1$ closer to
$D_1 = 70$ and assigns the larger shortfall to stream 2, while still respecting
both minimum guarantees ($G_1=40$, $G_2=30$) via the capacity-slack variable
$u$ that absorbs the remaining $100 - B_1 - B_2$ Mbps.

\end{document}
